%% ma: L12-B01-0010
%% lop: 12
%% chuong: 1
%% bai: 1
%% dang: 12.1.6
%% loai: TN4
%% muc: TH
%% dap_an: "C"
%% nguon: "(NBV)-ĐỀ SỐ 7-MỖI NGÀY 1 ĐỀ THI 2025.pdf (D07, MỖI NGÀY 1 ĐỀ - Đề số 7), Phần I Câu 1"
%% kien_thuc: "Bài 1 (cực trị của hàm phân thức), đạo hàm của thương (Lớp 11 Bài 32)"
%% trang_thai: cho-duyet
%% ly_do: "Dạng toán mới đề xuất 12.1.6 (Tìm cực trị của hàm phân thức bậc hai trên bậc nhất, xét khẳng định sai); chờ giáo viên duyệt dạng."
%% ghi_chu: "Đáp án gốc C đúng (kiểm bằng sympy)."
\begin{ex}
Cho hàm số $y=\dfrac{x^2+2x+8}{x-2}$. Khẳng định nào sau đây sai?
\choice{Hàm số đạt cực đại tại $x=-2$}{Hàm số đạt cực tiểu tại $x=6$}{\True Giá trị cực tiểu của hàm số là $y=6$}{Giá trị cực đại của hàm số là $y=-2$}
\loigiai{Tập xác định $\mathbb R\setminus\{2\}$. $y'=\dfrac{x^2-4x-12}{(x-2)^2}$, $y'=0\Leftrightarrow x=-2$ hoặc $x=6$. $y'$ đổi dấu từ $+$ sang $-$ tại $x=-2$ (cực đại, $y(-2)=-2$) và từ $-$ sang $+$ tại $x=6$ (cực tiểu, $y(6)=14$). Vậy khẳng định ở C sai: giá trị cực tiểu bằng $14$, không phải $6$.}
\end{ex}
